So far, we focused on $k = 1$, discarding a general integer value, and also we used only one specific comparison, $\leq$. As we stated, cardinality constraints 
are particular types of constraints used very often inside the SAT models, their formally definition is:
$$\sum_{1 \leq i \leq n} x_i \bowtie k$$
where:
\begin{itemize}
    \renewcommand{\labelitemi}{-}
    \item $\bowtie \rightarrow$ comparison, it belongs to the set $\{<, \leq, =, \neq, >, \geq\}$ 
\end{itemize}
However, let's take a look to the whole picture considering also other comparisons, in particular:
\begin{itemize}
    \renewcommand{\labelitemi}{-}
    \item \dots
\end{itemize}
As we can notice, every single comparison can be expressed as $\leq$ (less or equal) of a fixed integer number. Therefore, we could reduce our discussion to define 
the best way to encode the AtMostK(\dots) constraint, since it can be used for any other comparison. \vspace{7pt}

Another special case is:
$$\sum_{1 \leq i \leq n} x_i = k \rightarrow \sum_{1 \leq i \leq n} x_i \leq k \wedge \sum_{1 \leq i \leq n} x_i \geq k$$